\begin{tikzpicture}[font=\small]
  % ---------- greedy frontier ----------
  \begin{scope}
    \node[anchor=south, font=\small] at (4.0,3.2){Dijkstra：每次从边界上取出 $d$ 最小的点，\textbf{定死}它};

    \begin{scope}[on background layer]
      \fill[algteal, opacity=0.10] (-1.0,-1.6) rectangle (2.0,2.4);
    \end{scope}
    \node[algteal, font=\footnotesize, anchor=north] at (0.5,-1.7){已定死集合 $S$（$d$ 不再改变）};

    \node[done] (s) at (-0.3,0.4) {$s$};
    \node[done] (a) at (1.3,1.5) {$a$};
    \node[done] (b) at (1.3,-0.7) {$b$};
    \node[vis]  (c) at (3.4,1.5) {$c$};
    \node[vis]  (d) at (3.4,-0.7) {$d$};
    \node[vtx]  (e) at (5.5,0.4) {$e$};

    \draw[tedge] (s) -- node[above left, font=\footnotesize]{$2$} (a);
    \draw[tedge] (s) -- node[below left, font=\footnotesize]{$5$} (b);
    \draw[dedge] (a) -- node[above, font=\footnotesize]{$3$} (c);
    \draw[dedge] (a) -- node[right, font=\footnotesize, pos=0.3]{$1$} (d);
    \draw[dedge] (b) -- node[below, font=\footnotesize]{$4$} (d);
    \draw[dedge] (c) -- node[above right, font=\footnotesize]{$2$} (e);
    \draw[dedge] (d) -- node[below right, font=\footnotesize]{$6$} (e);

    \node[note, anchor=west, align=left] at (6.4,1.4)
      {\textcolor{algteal}{绿}：已定死\\
       \textcolor{algblue}{蓝}：在优先队列里（边界）\\
       白：还没碰到过};

    \node[note, anchor=north, align=center] at (3.4,-2.4)
      {正确性靠一条性质：\textbf{边权非负} $\Rightarrow$ 边界上 $d$ 最小的那个点，
       不可能再通过任何其他点绕出更短的路 ——\\ 因为绕出去只会让距离\textbf{不减}。};
  \end{scope}

  % ---------- why negative weights break it ----------
  \begin{scope}[yshift=-6.4cm, xshift=0.8cm]
    \node[anchor=south, font=\small] at (3.4,2.2){负权边为什么会让它出错};

    \node[vtx] (s2) at (0,0.8) {$s$};
    \node[bad] (u2) at (2.6,1.7) {$u$};
    \node[vtx] (v2) at (2.6,-0.2) {$v$};
    \draw[dedge] (s2) -- node[above left, font=\footnotesize]{$1$} (u2);
    \draw[dedge] (s2) -- node[below left, font=\footnotesize]{$4$} (v2);
    \draw[dedge, algred, line width=1.3pt] (v2) -- node[right, font=\footnotesize]{$-3$} (u2);

    \node[note, anchor=north west, align=left] at (-0.4,-0.9)
      {Dijkstra 会先取出 $u$（$d=1$ 最小）并\textbf{定死}它。\\
       但真实最短路是 $s\to v\to u$，代价 $4-3=1$… 这里恰好相等；\\
       把 $-3$ 换成 $-4$，真实答案就是 $0$，而 $u$ 早已被定死在 $1$ —— \textbf{错了}。\\[4pt]
       根因：非负性一旦失去，「当前最小」就不再意味着「最终最小」。\\
       这正是 Bellman-Ford 必须逐轮重扫所有边的理由。};
  \end{scope}

  % ---------- complexity ----------
  \node[note, anchor=north west, align=left] at (8.6,-5.2)
    {\textbf{复杂度取决于优先队列}\\[4pt]
     数组：$O(|V|^2+|E|)$\\
     二叉堆：$O((|V|+|E|)\log|V|)$\\
     斐波那契堆：$O(|E|+|V|\log|V|)$\\[5pt]
     稠密图用数组反而更快 ——\\
     $|E|\approx|V|^2$ 时 $\log$ 因子是纯亏。};
\end{tikzpicture}
